\egasection{§3\quad 泛超平面截面：几何不可约性与连通性}
\source{\texttt{EGAVaiello.pdf} 第12--13页}
\egas{3.1}{定理（Bertini--Zariski）}
假设 $\dim\overline{f(X)}\geq2$，且 $X$ 几何不可约。则泛超平面截面 $Y_\eta$ 也几何不可约。
\begin{proof}
设 $K/k$ 是 $X$ 的函数域，$n=\dim P$。在 $P$ 中选取仿射坐标 $T_1,\ldots,T_n$（选择无穷远超平面 $H^\infty$，使 $f(X)$ 不包含于其中），并在 $P^\vee$ 中取仿射坐标 $S_1,\ldots,S_n$。可见 $Y_\eta$ 的函数域 $L$ 可与整环
\[
 K[S_1,\ldots,S_n]\Big/\Bigl(\sum_{i=1}^n t_iS_i-1\Bigr)
\]
的分式域相等同，其中 $t_i\in K$ 是 $T_i$ 沿 $f:X\to P$ 的像。由于 $\dim\overline{f(X)}>0$，这些 $t_i$ 并非全在 $k$ 上代数，因而更不可能全为零；例如可设 $t_n\ne0$。于是立即得到
\[
 L=K(S_1,\ldots,S_{n-1}),
\]
这是纯超越扩张；元素 $S_n\in L$ 由方程 $\sum_it_iS_i-1=0$ 表成各 $S_i$（$1\leq i\leq n-1$）和 $t_i$（$1\leq i\leq n$）的函数。
\cmt{B2-P01}{定理只假设几何不可约，没有假设约化。这里的“函数域”应先在约化不可约预概形 $X_{\mathrm{red}}$ 上解释，或使用泛点的剩余域。这样补足不改变所研究的几何不可约性；不是把原假设偷偷加强为几何整。}
另一方面，$k'=k(\eta)$ 可与 $k(S_1,\ldots,S_n)$ 等同，而典范嵌入 $k'\to L$ 由把 $S_i$ 映到 $S_i$ 得到〔PB：核对这一点！〕。
\ednote{大概应是把 $S_i$ 映到它在 $L$ 中的等价类 $[S_i]$。}
换言之，作为 $L$ 的子扩张，$k'$ 由这些 $S_i$（$1\leq i\leq n$）生成；等价地，由 $S_1,\ldots,S_{n-1}$ 和
\[
 S_n=a_0+a_1S_1+\cdots+a_{n-1}S_{n-1},
 \qquad a_0=t_n^{-1},\quad a_i=-t_it_n^{-1}\ (1\leq i<n)
\]
生成。注意，各 $a_i$ 生成的域与各 $t_i$ 生成的域显然相同；它们在 $k$ 上的共同超越次数就是 $\overline{f(X)}$ 的维数。（附记：把这个双有理描述至少作为（2.1）的推论收入，应当是合适的。）所以（3.1）的证明归结为下面的引理。
\end{proof}

\egas{3.1.1}{引理（Zariski）}
〔参看本节末的英译编者注。〕设 $k$ 为域，$K/k$ 为有限型域扩张，$m\geq0$ 为整数，$a_i\in K$（$0\leq i\leq m$），且 $k(a_0,\ldots,a_m)$ 在 $k$ 上的超越次数至少为2。令
\[
 L=K(S_1,\ldots,S_m),\qquad
 k'=k\left(S_1,\ldots,S_m,a_0+\sum_{i=1}^m a_iS_i\right)\subset L,
\]
其中 $S_i$ 为不定元。如果 $K$ 是 $k$ 的准素扩张，则 $L$ 是 $k'$ 的准素扩张。
\ednote{primary extension 大概表示较小的域在较大的域中代数闭，或者准代数闭。参见 Jouanolou，定理3.6。}
\cmt{B2-H07}{这里保留英译编者的猜测，但它不能替代定义。所上传 EGA IV 第二部分（4.3.1），PDF第55页，明确说：$L/k$ 为准素扩张，当且仅当 $L$ 内包含的 $k$ 的最大可分代数扩张就是 $k$ 本身。它不要求一切代数元素都已在 $k$ 中；非平凡纯不可分扩张就说明二者不同。}

这一引理，或者同它如同兄弟般相似的引理，在文献中几乎随处可见。因此，我把选取一份证明出处的工作留给你；也就是说，我没有兴致凭自己的办法再找一个证明。
\cmt{B2-P02}{这句话正是上传笔记的内容；此处没有引理的证明。故（3.1）在本载体内只是归结到一个外部引理，不能宣称本节已经给出了自足的完整证明。这里不补造原文未写的证明。}

\egas{3.2}{推论}
假设 $f$ 非分歧，或者 $k$ 的特征为0，并且 $\dim\overline{f(X)}\geq2$。如果 $X$ 几何整，则 $Y_\eta$ 也几何整。
\begin{proof}
因为“几何整”等价于“几何不可约且可分”。
\end{proof}

\egas{3.3}{推论}
假设 $k$ 代数闭，对 $X$ 的每个不可约分支 $X_i$ 都有 $\dim\overline{f(X_i)}\geq2$。再设 $X$ 是 $\sigma$-连通的，其中 $\sigma$ 为满足 $\dim\overline{f(Z)}=0$ 的 $X$ 的闭子集 $Z$ 所组成的集合；这就是说，对于每个这样的 $Z$，$X-Z$ 都连通。在这些条件下，$Y_\eta$ 在 $k(\eta)$ 上几何连通。
\begin{proof}
按照一个应当与 Hartshorne 定理一起出现在第6段中的引理，\ednote{询问 A. G.。}
上述假设表示：任意两个不可约分支 $X',X''$，都能由不可约分支链
$X_0=X',\ldots,X_n=X''$ 联系起来，使相邻两项的交不属于 $\sigma$。因此它们的逆像 $X'_\eta,X''_\eta$ 由一条 $X_{i\eta}$ 链相连；由（3.1），各 $X_{i\eta}$ 在 $k(\eta)$ 上几何连通，由（2.3），相邻两项的交非空。因为 $Y_\eta$ 是所有这些 $X_{i\eta}$ 的并，故 $Y_\eta$ 在 $k(\eta)$ 上几何连通，证毕。
\cmt{B2-T06}{本段英译稿把链的首项写成“$X_0$ and $X'$”，又在末句写 $Y_\eta=X_\eta$。这里分别按同句定义明确排为 $X_0=X'$ 和 $Y_\eta=\bigcup_iX_{i\eta}$，记录这一记号整理，不将其另算为定理反例。}
\end{proof}
\noindent\textbf{英译编者关于（3.1.1）的注。}
应当与 Zariski 的\textit{论文集}（MIT Press）第1卷第174页、第2卷第304页比较；还可参看 Zariski--Samuel，Springer 的 GTM 版第1卷第196页、第2卷第230页，以及 Jouanolou 的 \textit{Théorèmes de Bertini et applications}，定理3.6和第6节。

\egasection{§4\quad 变动超平面截面的研究：“充分一般”的截面}
\source{\texttt{EGAVaiello.pdf} 第13--15页}
我们回到§1的一般情形，$S$ 为任意预概形；另假设 $X$ 在 $S$ 上有限表示。

一般说，若 $P(Z,k)$ 是域 $k$ 上代数预概形 $Z$ 的一种“可构造”性质，则满足 $P(Y_\xi,k(\xi))$ 的 $\xi\in P^\vee$ 组成可构造集合：只须注意 $Y_\xi$ 是有限表示态射 $Y\to P^\vee$ 在 $\xi$ 处的纤维，并应用第9段。\ednote{指 EGA IV。}
对于性质 $P(Z,M,k)$ 也有类似说明，其中 $Z,k$ 同前，$M$ 是 $Z$ 上的凝聚模；底本接着假设“$G$ 在 $X$ 上有限表示”，并写满足 $P(Y_\eta,G_\eta,k(\eta))$ 的 $\xi\in P^\vee$ 所成的集合可构造。
\cmt{B2-T07}{此处前后使用的载体、变动指标不一致。按§1，$F$ 在 $X$ 上，$G$ 在 $Y$ 上；而讨论随 $\xi$ 变化的纤维时，应写 $P(Y_\xi,G_\xi,k(\xi))$。保持底本措辞可看清这一问题，不能据此把一个固定泛纤维性质直接当成各纤维性质。}
另一方面，在前两节中，当 $S=\Spec k$ 时，我们已在若干情形给出了上述集合 $E$ 包含 $P^\vee$ 泛点的判据。由可构造性，$E$ 便包含非空开集：这正是从泛超平面截面的结论，过渡到“充分一般”超平面截面的相应结论。

还应注意，若 $S=\Spec k$，则有逆向表述：泛超平面截面具有性质 $P$，当且仅当 $\eta$ 的某个合适邻域内各 $Y_\xi$ 都具有该性质；只要求其中的闭点 $\xi$ 满足它，也已足够。若 $k$ 代数闭，这归结为考虑 $k$-有理点，也就是 $X$ 本身的超平面截面，而无须先扩张底域。这来自 $E$ 的可构造性以及 $P^\vee$ 是 Jacobson 概形的事实。下面列举几个适用的特殊情形。

\egas{4.2}{命题}
设 $G$ 是 $X$ 上的有限表示模。
\cmt{B2-T07续}{底本此处确实写 $G$，随后的(b)又写 $F$；这里仍标明此处原字母，后文按各自原文保留。原稿没有单独印出“4.1”条目标记，虽然以后用4.1指本节开头的讨论，译文不补造一个新的原编号。}
设 $P$ 是域 $K$ 上代数概形 $Z$ 的模 $M$ 的下列性质之一：
\begin{enumerate}[label=(\roman*)]
\item $\coprof(M)\leq n$；
\item $M$ 满足 $(S_k)$；
\item $M$ 是 Cohen--Macaulay 模；
\item $M$ 没有嵌入的相伴素循环分支；
\item $M$ 在 $K$ 上可分。
\end{enumerate}
在这些记号下，令 $E$ 为使 $G_\xi$ 满足性质 $P$ 的 $\xi\in P$ 的集合。则：(a) $E$ 是 $P^\vee$ 的可构造子集；(b) 若 $S=\Spec k$，$k$ 为域，且 $F$ 满足性质 $P$，再在(v)的情形假设 $k$ 的特征为0或者 $f:X\to P$ 非分歧，则 $E$ 包含 $P^\vee$ 的一个稠密开子集。
\cmt{B2-T07续}{这里的参数应在 $P^\vee$ 而不是 $P$ 中；(a)以及§1的参数空间定义已给出内部依据。}
\begin{proof}
(a)来自我们在 EGA IV 第9段中见到的事实：$P$ 是可构造性质。(b)则来自前两节的相应结果。
\end{proof}
\noindent\textbf{对(b)的补充念头。}更一般地，设 $Z$ 是 $X$ 上 $F$ 不满足 $P$ 的点集，只要求 $f(Z)$ 有限，即 $\dim\overline{f(Z)}\leq0$。

\egas{4.3}{命题}
设 $P$ 是域 $K$ 上代数预概形的下列性质之一：
\begin{enumerate}[label=(\roman*)]
\item $Z$ 在 $K$ 上光滑；
\item $Z$ 在 $K$ 上满足几何 $(R_k)$；
\item $Z$ 在 $K$ 上可分；
\item $Z$ 在 $K$ 上几何正规；
\item $Z$ 在 $K$ 上几何整；
\item $Z$ 在 $K$ 上几何不可约。
\end{enumerate}
令 $E$ 为使 $Y_\xi$ 满足 $P$ 的 $\xi\in P^\vee$ 的集合。则：(a) $E$ 为 $P^\vee$ 的可构造子集；(b) 假设 $S=\Spec k$，$k$ 为域；在(i)至(v)的情形，假设 $k$ 的特征为0，\emph{并且} $f:X\to P$ 非分歧。最后，在(v)、(vi)的情形假设 $\dim\overline{f(X)}\geq2$。若 $X$ 满足 $P$，则 $E$ 包含 $P^\vee$ 的一个稠密开子集。
\cmt{B2-H08}{本段英文使用 and，译文保留“并且”。前面（3.2）使用 or；把两个充分条件同时要求只是更强的假设，不能仅因其较强就断言这里的定理为假，也不应静默把 and 改成 or。}
\begin{proof}
证明与（4.2）相同。还应指出，在(i)至(v)的情形，对(b)只假设 $f(Z)$ 有限便足够，其中 $Z$ 是 $X$ 上性质 $P$ 不成立的点集。
\end{proof}

\egas{4.4}{推论}
考虑性质 $(C_m)$：“$\overline Z$ 不能被一个维数 $\leq m$ 的闭子集断开”，其中
$\overline Z=Z\times_K\overline K$，$\overline K$ 是 $K$ 的代数闭包。令 $E$ 为使 $Y_\xi$ 在 $k(\xi)$ 上满足 $(C_m)$ 的 $\xi\in P^\vee$ 的集合。则：(a) $E$ 可构造；(b) 假设 $S=\Spec k$，$k$ 为域，并且对 $X$ 的每个不可约分支 $X_i$，都有 $\dim\overline{f(X)}\geq2$。若 $X$ 在 $k$ 上满足 $(C_{m+1})$，则 $E$ 包含 $P$ 的一个稠密开子集 $U$。
\cmt{B2-E02}{底本此处还有 $f(X)$ 与 $f(X_i)$、$P$ 与 $P^\vee$ 的记号错置；但即便把这两个记号整理成预期形式，一般态射下的连通性断言仍不成立。报告给出 $X=V(uv)\subset\mathbf A^4$、$f=[1:u+v:(u+v)w]$ 的具体反例：$X$ 满足 $(C_1)$，而泛截面不满足 $(C_0)$。}

可构造性通过 AQT 得到。\ednote{AQT 是什么？询问 AG。}
这是 EGA IV 第9段里遗漏的一个事实，也许还来得及补上。(b)原则上来自（3.3），只是（3.3）的表述过于特殊，因而应作推广；证明基本不变。此外，（4.4）的陈述中有一个错误，它按现在的写法“在 $f$ 拟有限时不成立”；在一般情形中，应当不再考虑 $X$、$Y_\xi$ 中闭子集本身的维数，而考虑它们在 $P$、$H_\xi$ 中的像的维数。\textit{Redactor demerdetur.}
\ednote{拉丁文；待翻译。}
\cmt{B2-P03}{所上传英文确实写“not valid as such if $f$ is quasi-finite”，这一句与随后“在一般情形应看像维数”的说明关系可疑，可能缺少否定或限制词；未见法文手稿，不能替作者裁定。拉丁语短句也保留原字，暂不强作可靠译释。明确反例针对一般非拟有限态射，不宣称它同时推翻拟有限版本。}

\egasection{§5\quad Seidenberg 类型定理}
\source{\texttt{EGAVaiello.pdf} 第16--18页的相应部分}
\egap{5.1}
本节给出（4.1）所定义集合 $E$ 为开集的条件。这里研究的是对 $X$、分别对 $Y_\xi$ 为局部性质的 $P$，从而可以定义 $Y$ 中的点集 $U$：若 $y$ 在 $P^\vee$ 中的像为 $\xi$，则要求 $Y_\xi$ 在 $y$ 处满足 $P$，分别要求 $G_\xi$ 在 $y$ 处满足 $P$。我们先利用第12段给出 $U$ 为开集的判据。\ednote{须查明引用；很可能是 EGA IV。是的。}
如常有 $E=P^\vee-h(Y-U)$〔英译编者补字〕。因此，若 $U$ 为开集且 $X$ 在 $S$ 上固有，则因 $h$ 固有、从而闭，$E$ 也为开集。
\ednote{旁注：因为 $Y$ 在 $X$ 上固有，而 $P^\vee$ 在 $S$ 上分离。}

\egap{5.2}
如§1所见，$Y$ 在 $X\times_SP^\vee=X_{P^\vee}$ 中由
$\OO_X(1)\otimes_S\OO_{P^\vee}(1)$ 的截面 $\phi$ 的“除子”定义。对每个 $\xi\in P^\vee$，它诱导
\[
 \OO_X(1)\otimes_{k(s)}\OO_{P^\vee}(1)(\xi)
\]
的截面 $\phi_\xi$。底本说，这个层非典范地同构于
$\OO_X(1)\otimes_{k(s)}k(\xi)=\OO_{X_{k(s)}}(1)$；$Y_\xi$ 就是该截面的“除子”。
\cmt{B2-T08}{这里应把 $\OO_X(1)$ 先限制到 $X_s$，而扩张到的纤维是 $X_{k(\xi)}$，不是右端所印的 $X_{k(s)}$。这是载体及底域的记号问题。}
（附记：可逆模 $L$ 的截面 $\phi$ 的除子，定义为同态
$\phi^{-1}:L^{-1}\to\OO$ 的像理想所定义的闭子预概形。）若 $F$ 为 $X$ 上的模层，其在 $Y$ 上的逆像，即 $F\otimes_{\OO_S}\OO_{P^\vee}=F_{P^\vee}$ 在闭子概形 $Y\subset X_{P^\vee}$ 上的逆像，正是同态
\[
 \phi^{-1}\otimes\id_{F_{P^\vee}}:
 F_{P^\vee}(-1,-1)\longrightarrow F_{P^\vee}
\]
的余核；正如 Mike 所说，记号 $(-1,-1)$ 不言自明。
\ednote{我猜这里指 Mike Artin。——P. B.}
同样，$G_\xi$ 是相应同态 $F_{k(\xi)}(-1,-1)\to F_{k(\xi)}$ 的余核，其中 $\xi$ 为 $P$ 中的点；若 $\xi$ 不是 $P^\vee$ 中的普通点，而是一个取值于 $S$ 上的 $S'$ 的点，也有相应解释。\cmt{B2-T08续}{参数点所在的 $P$ 应为 $P^\vee$，同段下一句已经再次使用后者。}

一般地，若 $L$ 是某处的可逆模，$\phi$ 为定义子预概形 $V(\phi)$ 的截面，则对每个模 $F$，其在 $V(\phi)$ 上的逆像，按通常的滥用语言，与
$\id_F\otimes{}\,\cdots:F\otimes L^{-1}\to F$ 的余核等同。\ednote{底本此处有“??”，同态的另一张量因子缺失。}
若这个同态单射，则称 $\phi$ 是 $F$-正则的。底本接着说：“若选取 $F$ 与 $\OO_X$ 的一个同构，这在局部是可能的，使 $\phi$ 与 $\OO_X$ 的一个截面等同，那么这一术语与此前引入的术语相容。”
\cmt{B2-T08续}{可在局部平凡化的是可逆模 $L$，不是任意的 $F$。例如 $F=\OO_X^{\oplus2}$ 一般不可能同构于 $\OO_X$。将 $L$ 平凡化后，$\phi$ 给出函数，以此函数在 $F$ 上的乘法是否单射来定义 $F$-正则。$\phi^{-1}$ 在这里是诱导的对偶同态记号，并不表示一个处处存在的倒数函数。}

\egas{5.3}{命题}
使用上述记号，令 $U$ 为 $X_P$ 中满足下列条件的点 $x$ 的集合：若 $x$ 在 $P^\vee$ 中的像为 $\xi$，则 $\phi_\xi$ 在 $x$ 处是 $F_{k(\xi)}$-正则的。于是：
\begin{enumerate}[label=(\alph*)]
\item 若 $F$ 有限表示且相对 $S$ 平坦，则 $U$ 为开集，而 $G|_U$ 相对 $P^\vee$ 平坦。
\item 对每个 $s\in P^\vee$，若 $\eta$ 为 $P^\vee_s$ 的泛点，则 $U$ 包含 $X_{k(\eta)}$。
\end{enumerate}
\cmt{B2-T09}{底本在此写 $X_P$，但参数化的空间为 $X_{P^\vee}$；(b)中要谈 $P^\vee$ 在 $s$ 上的纤维，应有 $s\in S$。此外 $G$ 原定义在 $Y$ 上，写 $G|_U$ 时应理解为限制到 $Y\cap U$，或先沿闭浸入直接像延拓到 $X_{P^\vee}$。}
\begin{proof}
(a) 因为 $F_{P^\vee}$ 有限表示且相对 $P^\vee$ 平坦，结论来自11.3；若 $S$ 局部诺特，也来自 $0_{\mathrm{III}}\ldots$。\ednote{查明这一引文。指 EGA IV。}

(b) 可设 $S=\Spec k$。由第3段，$F_{k(\eta)}$ 的相伴循环是 $F$ 的相伴循环 $Z_i$ 的逆像。\ednote{字迹不清；询问 AG 或自行辨认——大概是 $\phi^{-1}$。}
若 $f(Z_i)$ 有限，则由（2.3），$Z_{i,k(\eta)}\cap Y=\varnothing$；否则，例如由（2.6），也有
$Z_{i,k(\eta)}\cap Y=Z_{i,k(\eta)}$（出于维数的理由；已经在（2.6）中用过的（2.3），无疑是更好的理由）。这说明 $\phi$ 不在任何 $Z_{i,k(\eta)}$ 上恒等于零，因而证明(b)。
\cmt{B2-T10}{上一句等号与下一句结论相反；这里应是 $Z_{i,k(\eta)}\cap Y\ne Z_{i,k(\eta)}$。等号意味着整个相伴循环被包含在零点集内，恰不能证明非零因子性。译文保留所上传英文的等号，并在此明示问题。}
\end{proof}

\egas{5.4}{推论}
令 $V$ 为使 $\phi_\xi$ 是 $F_{k(\xi)}$-正则的 $\xi\in P^\vee$ 的集合。若 $F$ 有限表示，则 $V$ 可构造，并包含 $P^\vee$ 在 $S$ 上各纤维的泛点。若另假设 $X$ 在 $S$ 上固有且 $F$ 在 $S$ 上平坦，则 $V$ 为开集。

\egas{5.5}{备注}
取位于 $s\in S$ 上方的 $\xi\in P^\vee$，并假设 $F_s$ 没有嵌入相伴循环。底本写道：“于是立即看到 $\xi\in V$（使用（5.4）的记号），这也就是说 $\Supp F_{k(\xi)}$ 的每个不可约分支都不位于 $H_\xi$ 上”；还说，在此判据中以 $k(\xi)$ 的任意扩张替换 $k(\xi)$，也成立，不过这一点稍不明显。
\cmt{B2-E03}{按字面，“没有嵌入相伴循环”不推出任意 $\xi\in V$。取 $F=i_*\OO_L$ 为 $\mathbf P^2$ 中直线 $L$ 的结构层，选 $H_\xi=L$，截面在 $F$ 上的乘法为零，因而不正则。合理的读法是把后半句作为 $\xi\in V$ 的充要条件，而不是从无嵌入循环自动推出它；详见报告。}
注意，刚才对 $F_s$ 所加的 $(S_1)$ 假设，在 $F_s$ 为 Cohen--Macaulay 时尤其成立，更不用说 $F$ 相对 $S$ 为 Cohen--Macaulay 的情形。笔记接着说，在这种情况下 $G_s$ 为 Cohen--Macaulay，因为局部上它由同样为 Cohen--Macaulay 的 $F_{k(s)}$ 模去 $\phi\cdot F_{k(s)}$ 得到，其中 $\phi$ 对 $F_{k(s)}$ 正则。上述说明同样应该、并且将会在前面作局部表述，以刻画 $U$ 的点而不是 $V$ 的点。
\cmt{B2-E03续}{这里使用模去正则元素的结论，仍须保留“截面正则”这一条件；不能只靠 Cohen--Macaulay 性就略去它。纤维下标 $s$ 与 $\xi$ 也应按具体基变换区分。}

现在利用12.1.1与12.1.4，得到：\ednote{询问 AG 这些引文——大概是 EGA IV。12.1.4核对不上。}
\egas{5.6}{定理}
假设 $F$ 有限表示且相对 $S$ 平坦。令 $P$ 为12.1.1中(i)至(viii)的性质之一，或者在假设 $F=\OO_X$ 时，为 EGA IV 的12.1.4中(i)至(iv)的性质之一〔英译编者补入卷号〕。以 $U_P$ 表示 $X_P$ 中的点 $x$ 的集合：若 $\xi$ 是其在 $P^\vee$ 中的像，则要求 $G_\xi$（分别地，$Y_\xi$）在 $x$ 处满足 $P$，且 $\phi_\xi$ 在 $x$ 处为 $F_{k(\xi)}$-正则。则 $U_P$ 为开集，且 $G|_{U_P}$ 相对 $S$ 平坦。
\begin{proof}
按定义，$U_P\subset U$（使用（5.3，a）的记号）；把第12段应用于
$U\to P^\vee$ 和 $F_{P^\vee}|_U$ 即得。
\cmt{B2-P04}{英译编者已经明确说引文核对不上，所列性质未在这里实际展开；证明最后写 $F_{P^\vee}|_U$，而结论涉及除子上的余核 $G$，也需与引用的原命题逐项对应。本次不把未确定的八项、四项性质凭常识补进正文。}
\end{proof}

\egas{5.7}{推论}
假设 $F$ 有限表示、相对 $S$ 平坦，并且 $\Supp F$ 在 $S$ 上固有（例如 $X$ 在 $S$ 上固有）。令 $V_P$ 为 $\xi\in P^\vee$ 的集合，使 $G_\xi$（分别地，底本此处写 $Y$）满足性质 $P$，并且“它是 $F_{k(\xi)}$-正则的”。在这些条件下，$V_P$ 为开集；而且在任何情况下，即不加平坦性或固有性假设时，它都是可构造的。
\cmt{B2-P04续}{“正则”的主语应为截面 $\phi_\xi$；谈该参数对应的纤维时应为 $Y_\xi$。此处仍受（5.6）未明确列出性质清单的限制，未据这一句话作超出来源的普遍推广。}

从表述的角度说，我觉得不能把（5.6）保留为只指向另一本中的条件编号，而应明确列出我们要用的性质(i)、(ii)、……。还应指出，也许可以在（5.1）中说，“几何正规”的情形（当然 $S=\Spec k$）是 Seidenberg 的结果。\ednote{或者意为“为稳妥起见”。}

\egasection{§6\quad 任意超平面截面的连通性}
\source{\texttt{EGAVaiello.pdf} 第18页}
这里把已知的泛超平面截面几何连通性判据（3.3），同 Zariski 连通性定理结合，以得到任意超平面截面的连通性结论。

\egas{6.1}{命题}
设 $S=\Spec k$，$k$ 为代数闭域〔假设 $X$ 在 $k$ 上固有〕；\ednote{这一处字迹不清。}
并设对 $X$ 的每个不可约分支 $X_i$，都有 $\dim\overline{f(X_i)}\geq2$。最后，假设 $X$ 不能被满足 $\dim\overline{f(Z)}\leq0$ 的闭子集 $Z$ 断开。在这些条件下，对每个 $\xi\in P^\vee$，$Y_\xi$ 都几何连通。
\cmt{B2-H09}{方括号中的固有性出自上传英译稿的补读，不能省去。报告给出删去它后特殊超平面截面不连通的例子；该例不是对保留固有性版本的反例。}
\begin{proof}
由于没有一个 $f(X_i)$ 为有限集，可见 $Y$ 的每个不可约分支 $Y_i$ 都占优于 $P^\vee$；另一方面，$Y\to P^\vee$ 固有。底本在括号中解释说：“如果 $Y$ 在 $k$ 上固有，因为它在 $X$ 上固有，而 $X$ 在 $k$ 上固有。”由（3.3），$Y\to P^\vee$ 的泛纤维 $Y_\eta$ 几何连通。最后，$P^\vee$ 正则，因而几何单支。应用15.6.3——Zariski 连通性定理的一个变体——就能推出 $Y\to P^\vee$ 的所有纤维几何连通，证毕。\ednote{指 EGA IV。}
\end{proof}
实际上，用同类证明，不难在（4.4）所示的方向上推广（6.1）。如果你不愿为这个练习费心，至少也应把它作为备注提一下。还应说清，（6.1）并不排除维数过大的超平面截面。从编排的角度看，把所有连通性问题，包括（3.3）、（4.4），集中在同一节中，也许更清楚。


